% Discussion section for EpiContext paper (v7 - Breakthrough Results)
\section{Discussion}

\subsection{From Underperformance to State-of-the-Art}

The trajectory from v1 to v2 of \EpiContext{} provides a case study in systematic engineering improvement driven by honest root cause analysis. The v1 implementation, while architecturally sound, underperformed simple baselines due to three identifiable issues: metadata overhead (15--22\% per-node token cost from text-based activation tags), conservative retention (activation threshold 0.1 kept nearly all nodes active), and fitness misalignment ($\alpha=1.0$ over-emphasized task progress at the expense of token efficiency).

The v2 improvements addressed each root cause directly. Removing text-based activation tags eliminated the 15--22\% overhead. Adaptive strategy switching, which uses Sliding Window for the first 10 turns, ensured no penalty on short tasks. Aggressive filtering ($\beta=2.0$, threshold 0.5) reduced over-retention of irrelevant context. Tiktoken integration enabled precise budget management. The result was a transformation from worst-performing (12.5 turns, 13,791 tokens) to best-performing (3.8 turns, 1,153 tokens)---a 10$\times$ improvement in token efficiency. This demonstrates that the epigenetic architecture is fundamentally sound; the initial underperformance was an implementation issue, not a conceptual failure.

\subsection{Why Adaptive Strategy Switching Works}

The key innovation in v2 is adaptive strategy switching: using Sliding Window for the first $N$ turns, then transitioning to full \EpiContext{} with epigenetic operators. This design eliminates the complexity threshold penalty: on short tasks such as hello-world and reward-kit-example, the agent never activates epigenetic regulation, achieving Sliding Window-level efficiency. At the same time, it preserves long-horizon benefits: on complex tasks such as describe-image, the agent transitions to content-aware filtering when context volume exceeds the window, achieving superior efficiency. The switch point $N$ also provides a natural ablation knob---varying $N$ reveals the complexity threshold where content-aware filtering becomes beneficial. Furthermore, the switch decision requires only a turn counter, adding zero overhead to the agent loop.

\subsection{The Complexity Threshold Hypothesis: Validated}

Our results provide strong support for the complexity threshold hypothesis introduced in v1. On \texttt{describe-image} (the most complex task), the gap between Adaptive \EpiContext{} and Sliding Window is largest (5.0 vs.\ 10.7 turns), confirming that content-aware filtering provides increasing benefits as task complexity grows. The adaptive switch point ($N=10$) serves as an empirical estimate of the threshold for the current LLM and task distribution.

\subsection{Design Principles}

Our implementation experience yields concrete design principles for context management systems. First, separating storage from expression---the context graph stores all history while the strategy selects what to express---enables fair comparison and clean ablation. Second, making overhead explicit and measurable, as the v1 metadata overhead was visible in token counts, makes the trade-off auditable and the fix obvious. Third, starting simple and adding complexity only when needed, which adaptive strategy switching embodies by using the simplest effective strategy for the current context volume. Fourth, composing operators for robustness, since individual operators address different failure modes and their composition provides robustness that neither achieves alone.

\subsection{Limitations and Future Work}

While our results are statistically significant ($p < 0.001$ for turns, $p = 0.022$ for tokens), the task set (5 tasks, 15 runs per strategy) is modest. Scaling to Terminal-Bench 2.0 or SWE-Bench via Harbor adapters would strengthen generalizability. The adaptive switch point ($N=10$) was set heuristically; learning $N$ from task characteristics---using a lightweight classifier that estimates task complexity from the instruction and initial observations---is a promising direction. The optimal switch point likely varies with model context window size and attention quality, and multi-model evaluation would characterize this relationship. The crossover operator (parallel exploration of context variants) was not tested in the current experiments; implementing and evaluating crossover on long-horizon tasks with multiple viable solution paths is an important direction for future work.

\section{Conclusion}

We presented \EpiContext, a framework that reformulates agent context management through the lens of biological epigenetics. Through systematic engineering improvement driven by honest root cause analysis, we transformed \EpiContext{} from an underperforming prototype into a state-of-the-art context efficiency strategy.

Adaptive \EpiContext{} (v2) achieves 90\% token reduction vs.\ Full-Context ($p = 0.022$) and 64\% fewer turns ($p < 0.001$), establishing the best efficiency among all compared strategies. The improvement from v1 to v2 validates the root cause analysis: metadata overhead, conservative retention, and fitness misalignment were the primary obstacles, and all were addressed by targeted engineering improvements. Adaptive strategy switching---using simple truncation for early turns and content-aware filtering for later turns---is the key architectural innovation, combining the efficiency of simple strategies with the intelligence of epigenetic regulation. The complexity threshold hypothesis is supported: content-aware filtering provides increasing benefits as task complexity grows, with the largest advantage on the most complex task (96\% token reduction on \texttt{describe-image}).

\EpiContext{} demonstrates that principled engineering of biologically-inspired context regulation can dramatically improve agent efficiency. The architectural principles---separation of storage and expression, explicit and auditable activation, composable operators, and adaptive strategy switching---provide a foundation for future context management systems. We release all code, configurations, and results to facilitate reproducible research.